\section{Related Work}
\label{sec:related}
Recent advances in Artificial Intelligence have significantly expanded the role of urban analytics in Smart City ecosystems. However, most existing studies focus on metropolitan environments, while comparatively limited attention has been devoted to small historic destinations.

% Pedoni
Pedestrian and crowd-flow forecasting are commonly framed as spatio-temporal
prediction problems in which recent flow history is combined with calendar,
weather, and spatial information. Citywide crowd-flow models often assume dense
grids, graphs, or transport networks. ST-ResNet, for example, introduced a
deep residual architecture for citywide crowd-flow prediction on spatial grids
\cite{Zhang2017STResNet}, while graph-based recurrent models such as DCRNN
represent flows over transportation networks \cite{Li2018DCRNN}. Pedestrian
volume forecasting has also been studied at high spatio-temporal granularity
using enhanced learning models \cite{Jiang2022PedestrianGranularity}.

% Turism forecasting
Tourism forecasting has long combined time-series models, econometric variables,
and increasingly diverse digital traces. Reviews emphasize that forecasting
performance is context-dependent and that no single model class dominates all
tourism-demand problems \cite{Song2019TourismDemandReview}. Recent work has
explored deep learning (DL) \cite{Law2019TourismDemandDeep} and multi-source big data
for demand forecasting \cite{Li2020TourismBigData}.

The Dozza setting differs from these large-scale urban setups. The available
pedestrian data is not a dense sensor graph but a small set of access-point
cameras. This makes a tabular approach with lagged histories, calendar variables,
weather covariates, and mobile-network aggregates a defensible baseline before
moving toward graph or deep sequence architectures. %It also makes baseline comparison especially important: at a one-hour horizon, recent local counts can already explain a substantial fraction of the next value.

Events can strongly affect local visitor dynamics. Mobile-phone data have been
used for tourism event analytics \cite{Leng2021TourismEventAnalytics} and
event-aware architectures have been proposed for mobility nowcasting
\cite{Wang2022EASTNet}. Our case study uses a lighter tabular alternative:
event records are transformed into hourly intensity, distance, city, category,
lag, and rolling features. \red{They are evaluated under the same chronological
forecast-origin protocol as the other covariates: for each horizon, model inputs
are built only from information that would be available at the prediction time,
while future target values and same-hour target-family variables are excluded.}
This is intentionally more modest
than a dedicated event-aware neural architecture, but it is appropriate for a
small local dataset and makes ablation against non-event feature groups direct.

Mobile and passive mobility data is particularly relevant when official counts
are sparse or delayed. Passive mobile data have been used to study seasonal
tourism mobilities \cite{Zaragozi2021PassiveMobileTourism}, and mobility-derived
variables can improve tourism demand models in some contexts
\cite{Li2022MobilityTourismForecasting}. Sparse geolocation signals have also
been explored for tourism-flow prediction with weather and holiday covariates
\cite{Lemmel2023SparseGeolocation}. The present work is local and short-horizon:
it does not forecast regional arrivals, but hourly flows and mobile-network
presence indicators in a specific historic borough. \red{Lemmel et
al.~\cite{Lemmel2022TourismFlowLimitedData} provide a closely related
limited-data tourism-flow benchmark: they compare deep-learning architectures
with ARIMA for visitor-flow prediction and show that neural models can exploit
additional contextual inputs while retaining fast inference. Their study is
primarily a model-family comparison on tourism-flow observations, whereas the
present work uses a broader Smart City perspective by integrating pedestrian
sensors, mobile-network aggregates, weather, and event records within a unified
AI-enabled decision-support framework.} %Furthermore, our analysis extends beyond model comparison by considering multiple target families, multi-horizon forecasting, feature-selection strategies, ablation studies, and explainability analyses aimed at supporting operational decision making in smart tourism destinations.

% Eventi


%Finally, because tabular ensemble models can be difficult to interpret, feature importance and local attribution methods are relevant. We use model importance, permutation importance, ablation, and SHAP summaries. SHAP provides a unified additive feature-attribution framework \cite{Lundberg2017SHAP}; here it is used as a diagnostic tool, not as causal evidence.
